%% language-de-test — German multilingual paper (acmart `language=german`).
%% German main language with an English translated title/keywords. Verifies the
%% German fixed strings — \keywordsname ("Zusätzliche Schlagwörter und Phrasen"),
%% \proofname ("Beweis"), \acksname ("Danksagungen"), \tablename ("Tabelle") —
%% while the figure name stays "Fig." (acmart sets it globally, acmart.dtx:4199).
%% Twin of language-de-test.typ.
\documentclass[acmsmall,language=german]{acmart}

\acmJournal{JACM}
\acmVolume{37}
\acmNumber{4}
\acmArticle{111}
\acmYear{2018}
\acmMonth{8}
\acmDOI{XXXXXXX.XXXXXXX}
\setcopyright{acmlicensed}
\copyrightyear{2018}

\begin{document}

\title{Eine Notiz über Berechnungskomplexität}
\translatedtitle{english}{A note on computational complexity}

\author{Hans Müller}
\email{mueller@example.de}
\affiliation{%
  \institution{Forschungsinstitut}
  \city{Berlin}
  \country{Germany}}

\renewcommand{\shortauthors}{Müller}

\begin{abstract}
  Eine kurze Notiz auf Deutsch, um den mehrsprachigen Modus der Klasse acmart mit
  übersetzten Zeichenketten zu prüfen.
\end{abstract}

\keywords{Komplexität, Algorithmen, Berechnung}
\translatedkeywords{english}{complexity, algorithms, computation}

\maketitle

\section{Einleitung}
Dieses Dokument prüft, dass die lokalisierten Zeichenketten zwischen der
LaTeX-Klasse und der Typst-Portierung übereinstimmen.

\begin{figure}[h]
\centering
\rule{4cm}{2cm}
\caption{Eine Abbildung; das Etikett bleibt englisch.}
\end{figure}

\begin{table}[h]
\caption{Eine Tabelle mit lokalisiertem Etikett.}
\centering
\begin{tabular}{ll}
\toprule
A & B\\
\midrule
1 & 2\\
\bottomrule
\end{tabular}
\end{table}

\begin{proof}
Ein kurzer Beweis, der mit einem QED-Quadrat endet.
\end{proof}

\begin{acks}
Wir danken den Gutachtern.
\end{acks}

\end{document}
